\section{Scope and Nonclaims}\label{sec:scope-nonclaims}

The HodgeCY II claims are deliberately narrower than the surrounding
geometric questions.  The verified statements are source-fidelity census
results, the \(84/84a\) source witness, exact degree-112 block-scheme
certification, the Hilbert--Burch structural explanation for the verified
block schemes, and the resulting source-versus-block-evaluation
non-determination statement.

The manuscript does not claim any of the following:

\begin{enumerate}
\item a complete geometric analysis of all \(114\) nontrivial source-level
  pairs or sets;
\item ordinary-node promotion for the \(84/84a\) block schemes;
\item equality between the verified block schemes and full saturated Jacobian
  node schemes;
\item an unconditional classical nodal-defect theorem;
\item a natural source-to-evaluation or source-to-node chain map;
\item an integral evaluation-lattice theorem;
\item a source-to-vanishing-cycle map;
\item an LMHS, MHM, or Hodge-atom identification;
\item a reverse implication from block-evaluation data to source assembly;
\item projective equivalence or inequivalence solely from matching or differing
  computed signatures.
\end{enumerate}

The status of each item should be read as open unless the final evidence matrix
states otherwise.  In particular, the equality \(h^1({\bf P}^3,\mathcal
I_B(8))=7\) is a theorem for the verified block schemes, not a promotion of
the classical nodal defect.
